% Part: computability
% Chapter: recursive-functions
% Section: composition

\documentclass[../../../include/open-logic-section]{subfiles}

\begin{document}

\olfileid[pt]{cmp}{rec}{com}
\olsection{Composição}

Se $f$ e $g$ são duas funções unárias de números naturais, podemos
compô-las: $h(x) = f(g(x))$. A nova função~$h(x)$ é então definida por
\emph{composição} a partir das funções $f$ e~$g$. Queremos generalizar isso
para funções de mais de um argumento.

Uma maneira de fazê-lo é a seguinte: suponha que $f$ seja uma função de
$k$ argumentos e que $g_0$, \dots, $g_{k-1}$ sejam $k$ funções, todas de
$n$ argumentos. Podemos então definir uma nova função de $n$ argumentos $h$
como segue:
\[
h(x_0, \dots, x_{n-1}) =
f(g_0(x_0, \dots, x_{n-1}), \dots, g_{k-1}(x_0, \dots, x_{n-1}))
\]
Se $f$ e todas as $g_i$ são computáveis, $h$ também é: para calcular
$h(x_0, \dots, x_{n-1})$, calculamos primeiro os valores
$y_i = g_i(x_0, \dots, x_{n-1})$ para cada $i = 0$, \dots,~$k-1$. Em
seguida, fornecemos esses valores a $f$ para calcular
$h(x_0, \dots, x_{k-1}) = f(y_0, \dots, y_{k-1})$.

Isso talvez pareça uma caracterização restritiva demais do que acontece
quando calculamos uma nova função usando funções já disponíveis. Para começar,
às vezes não usamos todos os argumentos de uma função, como quando definimos
$g(x, y, z) = \Succ(z)$ para empregá-la na definição por recursão primitiva
de~$\Add$. Suponha que possamos usar as seguintes funções:
\[
\Proj{n}{i}(x_0, \dots, x_{n-1}) = x_i
\]
As funções~$\Proj{k}{i}$ são chamadas de funções de \emph{projeção}:
$\Proj{n}{i}$ é uma função de $n$ argumentos. Nesse caso, $g$ pode ser
definida por
\[
g(x, y, z) = \Succ(\Proj{3}{2}(x, y, z)).
\]
Aqui o papel de $f$ é desempenhado pela função de $1$ argumento $\Succ$, de
modo que $k=1$. Temos ainda uma função de $3$ argumentos $\Proj{3}{2}$ que
desempenha o papel de $g_0$. O resultado é uma função de $3$ argumentos que devolve o
sucessor do terceiro argumento.

As funções de projeção também nos permitem definir novas funções reordenando
ou identificando argumentos. Por exemplo, a função $h(x) = \Add(x, x)$ pode
ser definida por
\[
h(x_0) = \Add(\Proj{1}{0}(x_0),\Proj{1}{0}(x_0)).
\]
Aqui $k=2$, $n=1$, o papel de $f(y_0,y_1)$ é desempenhado por $\Add$, e os
papéis de $g_0(x_0)$ e $g_1(x_0)$ são ambos desempenhados por
$\Proj{1}{0}(x_0)$, a função de projeção unária (isto é, a função identidade).

Se $f(y_0, y_1)$ é uma função que já temos, podemos definir a função
$h(x_0, x_1) = f(x_1, x_0)$ por
\[
h(x_0, x_1) = f(\Proj{2}{1}(x_0, x_1),\Proj{2}{0}(x_0, x_1)).
\]
Aqui $k=2$, $n = 2$, e os papéis de $g_0$ e $g_1$ são desempenhados,
respectivamente, por $\Proj{2}{1}$ e~$\Proj{2}{0}$.

Também pode parecer problemático exigir que $g_0$, \dots,~$g_{k-1}$ tenham
todas a mesma aridade~$n$. (Lembre-se de que a \emph{aridade} de uma função é
seu número de argumentos; uma função de $n$ argumentos tem aridade~$n$.) Mas
a inclusão das funções de projeção proporciona a flexibilidade desejada. Por
exemplo, suponha que $f$ e~$g$ sejam funções de $3$ argumentos e que $h$ seja
a função de $2$ argumentos definida por
\[
h(x,y) = f(x,g(x,x,y),y).
\]
A definição de~$h$ pode ser reescrita com as funções de projeção como
\[
h(x,y) = f(\Proj{2}{0}(x,y), g(\Proj{2}{0}(x,y), \Proj{2}{0}(x,y),
\Proj{2}{1}(x,y)), \Proj{2}{1}(x,y)).
\]
Então $h$ é a composição de $f$ com $\Proj{2}{0}$, $l$ e
$\Proj{2}{1}$, onde
\[
l(x,y) = g(\Proj{2}{0}(x,y),\Proj{2}{0}(x,y),\Proj{2}{1}(x,y)),
\]
isto é, $l$ é a composição de $g$ com $\Proj{2}{0}$, $\Proj{2}{0}$
e~$\Proj{2}{1}$.

\end{document}
